\section*{第 \ref{ch:b1:non-inertial-frames} 章——非惯性参考系；地球上的动力学}

\begin{solution}{exo:b1:non-inertial-frames:1}
$\tan\alpha = a/g = 0.255$，$\alpha = 14^\circ$（摆锤向后偏）；
$T = m\sqrt{g^2 + a^2} = 1.03\,mg$。刹车时偏角与张力相同，摆锤
向前倾斜。
\end{solution}

\begin{solution}{exo:b1:non-inertial-frames:2}
$\Omega = 2\pi/4.0 = \qty{1.57}{rad/s}$；$\Omega^2r = \qty{7.4}{m/s^2} = 0.75g$，
在转动\omterm{def:b1:point-kinematics:frame}{参考系}中沿外向；座位（摩擦与
约束）提供真实的向内力。
\end{solution}

\begin{solution}{exo:b1:non-inertial-frames:3}
$\Omega = 2\pi/86164 = \qty{7.29e-5}{rad/s}$；$\Omega^2R = \qty{0.034}{m/s^2}$，
相当于 $0.35\%$ 的 $g$。要使物体飘离地面需 $\Omega^2R = g$：$T = 2\pi\sqrt{R/g} =
\qty{84}{min}$——比现在快十七倍。
\end{solution}

\begin{solution}{exo:b1:non-inertial-frames:4}
$a_C = 2\Omega v\sin\lambda = 2 \times 7.29\times10^{-5} \times 27.8 \times 0.707
= \qty{2.9e-3}{m/s^2}$；$F = \qty{1.1}{kN}$。向北行驶时被推向东侧，
压向右侧（东侧）钢轨；向东行驶时被推向南，
同样压向右侧钢轨。
\end{solution}

\begin{solution}{exo:b1:non-inertial-frames:5}
$z = \Omega^2r^2/2g$：在 \qty{2}{rad/s} 时为 \qty{2.0}{mm}，在
\qty{10}{rad/s} 时为 \qty{5.1}{cm}。抛物面能将平行光束会聚于一点：
旋转的水银是一面完美、廉价且能自行抛光的抛物面镜
（且无法倾斜）。
\end{solution}

\begin{solution}{exo:b1:non-inertial-frames:6}
离心加速度 $\Omega^2R\cos\lambda$ 背离
转轴；它垂直于当地竖直方向的分量为 $\Omega^2R\cos\lambda
\sin\lambda$，它使 $\vect g$ 偏转 $\delta \approx \Omega^2R\sin\lambda\cos\lambda/g
= 0.034 \times 0.5/9.81 = \qty{1.7e-3}{rad} = 6'$。赤道处 $g$ 减小
$\Omega^2R = \qty{0.034}{m/s^2}$，源于自转（地球扁率还会使其进一步减小）。
\end{solution}

\begin{solution}{exo:b1:non-inertial-frames:7}
$\ddot x_E = 2\Omega gt\cos\lambda$: $x_E = \tfrac13\Omega g\cos\lambda\,t^3$; $t =
\sqrt{2h/g} = \qty{4.5}{s}$: $x_E = \tfrac13 \times 7.29\times10^{-5} \times 9.81
\times 0.707 \times 92 = \qty{1.6}{cm}$。向东偏：竖井顶部
比底部向东运动得更快（半径更大），石块保持
这一多出的向东速度——或者在转动\omterm{def:b1:point-kinematics:frame}{参考系}中，竖直向下运动的物体所受
\omterm{thm:b1:non-inertial-frames:dynamics}{科里奥利力}指向东方。
\end{solution}

\begin{solution}{exo:b1:non-inertial-frames:8}
$T_F = 23.93/\sin\lambda$：巴黎 \qty{31.8}{h}；极点 \qty{23.9}{h}；赤道
无穷大（无进动）。巴黎：$360^\circ/31.8 = 11.3^\circ$ 每小时。
\end{solution}

\begin{solution}{exo:b1:non-inertial-frames:9}
$a_C = 2\Omega v\sin\lambda = \qty{2.1e-3}{m/s^2}$，指向右侧。空气
从四面八方汇聚，向右偏转，因而盘旋内卷，
方向为逆时针（从上方看）：这是北半球的气旋。水槽
在数秒内排空，龙卷风在数分钟内形成：科里奥利偏转在
这样的时间尺度内，与任何初始旋涡相比都可以忽略。
\end{solution}

\begin{solution}{exo:b1:non-inertial-frames:10}
在汽车\omterm{def:b1:point-kinematics:frame}{参考系}中，视重力 $\vect g - \vect a$ 向后倾斜。
空气较重，向后“落”；气球比其排开的空气轻，
沿视竖直方向受到浮力托举——它向\emph{前}倾斜，
与铅垂线的方向相反。
\end{solution}

\begin{solution}{exo:b1:non-inertial-frames:11}
$\Omega^2R = g$: $\Omega = \qty{0.31}{rad/s}$, $T = \qty{20}{s}$. $a_C = 2\Omega v
= \qty{0.94}{m/s^2}$，约为 $0.1g$：顺着转动方向行走时，他感觉
体重增加了 $10\%$（\omterm{thm:b1:non-inertial-frames:dynamics}{科里奥利力}指向外侧），逆着转动方向则
感觉体重减轻。落下的小球：在\omterm{thm:b1:newton-dynamics:laws}{惯性参考系}中，它沿直线运动，而
地面向上弯曲；在下落时间 $\sqrt{2h/g} = \qty{0.64}{s}$ 内，
偏移量的\omterm{def:b1:units-dimensions:oom}{数量级}为 $\tfrac13\Omega\,g\,t^3 \cdot 2 \approx 2 \times 0.31 \times
9.81 \times 0.26/3 \approx \qty{0.5}{m}$——即几分米，与
转动方向相反。
\end{solution}

\begin{solution}{exo:b1:non-inertial-frames:12}
$GM/(d - r)^2 - GM/d^2 \approx 2GMr/d^3$（对 $r/d$ 取一阶近似）。月球：
$2 \times 6.67\times10^{-11} \times 7.35\times10^{22} \times 6.37\times10^6/
(3.84\times10^8)^3 = \qty{1.1e-6}{m/s^2}$；太阳：$2 \times 1.33\times10^{20}
\times 6.37\times10^6/(1.50\times10^{11})^3 = \qty{5.0e-7}{m/s^2}$；比值
$0.46$：太阳的\omterm{rem:b1:non-inertial-frames:tides}{潮汐}几乎是月球的一半，尽管其
引力大了 $180$ 倍——\omterm{rem:b1:non-inertial-frames:tides}{潮汐}与 $M/d^3$ 成正比。
\end{solution}

\begin{solution}{pb:b1:non-inertial-frames:1}
\textbf{1.} $T = 2\pi\sqrt{67/9.81} = \qty{16.4}{s}$; $x_0/\ell = 0.045$
($2.6^\circ$)：很小。

\textbf{2.} $v_{\max} = x_0\omega_0 = 3.0 \times 0.383 = \qty{1.15}{m/s}$；高度
$x_0^2/2\ell = \qty{6.7}{cm}$.

\textbf{3.} $T = mg + mv_{\max}^2/\ell = mg(1 + 0.002)$: $1.002\,mg$.

\textbf{4.} $\tfrac12\rho C_xSv_{\max}^2 = 0.5 \times 1.2 \times 0.05 \times 1.3 =
\qty{0.04}{N}$，与 $m\omega_0^2x_0 = \qty{12}{N}$ 相比小三百倍，
但足以在数小时内使摆锤停止。

\textbf{5.} 长线给出长周期，因而摆锤运动缓慢，
阻力的相对作用较小，并使每次摆动的进动清晰可见；
重摆锤储存了大量能量供阻力消耗——可摆动数小时。

\textbf{6.} 重力（竖直方向）与张力（位于线与竖直方向所确定的平面内，
即摆动平面内）：没有任何一个力具有
垂直于该平面的分量。

\textbf{7.} 摆面相对恒星固定，冰面在其下方以 $\Omega$ 向东转动，
于是观察者看到摆面顺时针转动，
经一个恒星日 \qty{23.93}{h} 后回到起点。

\textbf{8.} 每个半摆动描出一条略微转过角度的弦，构成花瓣形图样；
$86164/8.2 \approx 10500$ 个半摆动，每个再转过 $0.034^\circ$。

\textbf{9.} $\vect\Omega = \Omega(0, \cos\lambda, \sin\lambda)$.

\textbf{10.} 取 $\vect v' = (\dot x, \dot y, 0)$：
$\vect\Omega\wedge\vect v' = \Omega(-\sin\lambda\,\dot y,\ \sin\lambda\,\dot x,\ -\cos\lambda\,\dot x)$,
于是 $-2m\vect\Omega\wedge\vect v'$ 的水平分量为
$2m\Omega\sin\lambda\,(\dot y, -\dot x)$：只含 $\Omega\sin\lambda$。

\textbf{11.} $\ddot x = -\omega_0^2x + 2\omega_F\dot y$, $\ddot y = -\omega_0^2y -
2\omega_F\dot x$; $\omega_F = 7.29\times10^{-5} \times 0.753 = \qty{5.5e-5}{rad/s}
\ll \omega_0 = \qty{0.38}{rad/s}$.

\textbf{12.} $\ddot u = -\omega_0^2u + 2\omega_F(\dot y - j\dot x) = -\omega_0^2u -
2j\omega_F\dot u$.

\textbf{13.} $-\alpha^2 - 2\omega_F\alpha + \omega_0^2 = 0$: $\alpha = -\omega_F \pm
\sqrt{\omega_F^2 + \omega_0^2} \approx -\omega_F \pm \omega_0$.

\textbf{14.} 括号内是一个平面振动（过原点的
固定直线，对实数 $A = B$ 而言）；因子 $\eu^{-j\omega_Ft}$ 使整个
图样以 $\omega_F$ 顺时针转动，摆面顺时针进动，速率为
$\Omega\sin\lambda$。

\textbf{15.} $T_F = 2\pi/\omega_F = \qty{1.14e5}{s} = \qty{31.8}{h}$; $11.3^\circ$
每小时。

\textbf{16.} $\sin\lambda = 0$：$\vect\Omega$ 的局部竖直分量
为零；水平运动时，水平\omterm{thm:b1:non-inertial-frames:dynamics}{科里奥利力}为零。

\textbf{17.} $2m\omega_Fv_{\max} = 2 \times 28 \times 5.5\times10^{-5} \times 1.15 =
\qty{3.5}{mN}$；重力 \qty{275}{N}（其 $10^{-5}$），回复力
$m\omega_0^2x_0 = \qty{12}{N}$：小三千倍。

\textbf{18.} $\omega_FT/2 = 5.5\times10^{-5} \times 8.2 = \qty{4.5e-4}{rad}$:
$x_0 \times 4.5\times10^{-4} = \qty{1.4}{mm}$，每个半摆动。

\textbf{19.} 圆环上相距一厘米的钉子每隔 $\sim 7$
个半摆动便倒下一根，约一分钟。

\textbf{20.} 用手推一下会给出横向速度：摆锤随后
描出一个椭圆，椭圆本身的缓慢进动（有限\omterm{def:b1:wave-propagation:signal}{振幅}的效应）
掩盖了傅科进动。

\textbf{21.} 沿径向的冲击会沿摆动方向增加能量，却不增加
横向速度，因此摆面不受影响；而横向分量
会使摆进入椭圆运动，从而伪造或掩盖
进动。

\textbf{22.} 进动速率量度的是 $\Omega\sin\lambda$，即
地球自转的竖直分量；由天文学已知 $\Omega$，则
进动速率给出纬度（反之亦然）。

\textbf{23.} $\sin\lambda = 23.93/48 = 0.499$: $\lambda = 30^\circ$.

\textbf{24.} 其角\omterm{def:b1:newton-dynamics:conservation}{动量守恒}，因此其轴
相对恒星保持固定，而地球在其下方转动：陀螺罗经，
不借助磁铁即可找到真北。

\textbf{25.} \omterm{thm:b1:non-inertial-frames:dynamics}{科里奥利力} $-2m\vect\Omega\wedge\vect v'$；其水平
分量随 $\sin\lambda$ 增大；转一周的时间为 $23.93\,\text{h}/\sin\lambda$，
\qty{31.8}{h}，在巴黎。
\end{solution}
